% answers_short.tex —— 1963 年·全国卷 · 简短答案速查
% 由 tools/make_short_answers.py 自 answers.tex 的 \gksolhead/\gkanswer 提取，
% 并与 1963/1963_数学真题_答案要点.md 逐题比对后生成，勿手改（改源 MD 后重跑）。

\begin{gknotice}
\textbf{使用说明}：本节只给最终答案便于快速核对，完整推导见解析版。答案为编写者独立推导，\textbf{非官方标准答案}。
\end{gknotice}

\gkpart{小题}
\begin{itemize}[leftmargin=2.6em]
  \item[\textbf{1.}] $-3-2\sqrt2$。
  \item[\textbf{2.}] (1) 模数为 $2$；辐角主值为 $60^\circ$，一般辐角为 $360^\circ k+60^\circ$（$k\in\mathbb Z$）；(2) $-\sqrt3-i$
  \item[\textbf{3.}] $CD=\dfrac25$。
  \item[\textbf{4.}] 见解析版。
  \item[\textbf{5.}] $23.28^{-1.1}=0.03135$。
\end{itemize}

\gkpart{大题}
\begin{itemize}[leftmargin=2.6em]
  \item[\textbf{6.}] $x=n\pi\pm\dfrac\pi4$，或 $x=2n\pi-\dfrac\pi6$，或 $x=(2n+1)\pi+\dfrac\pi6$，其中 $n\in\mathbb Z$。
  \item[\textbf{7.}] (1) $720$ 个；(2) $240$ 个；(3) $120$ 个
  \item[\textbf{8.}] $(x,y)=(1,-2)$ 或 $\left(\dfrac{3\sqrt2}{2},\dfrac{\sqrt2}{2}\right)$。
  \item[\textbf{9.}] 见解析版。
  \item[\textbf{10.}] (1) $l+r=\dfrac{2}{\tan\theta(1-\tan^2\theta)}$；(2) $T=\dfrac{2\pi}{\tan^2\theta(1-\tan^2\theta)}$；(3) $\theta=\arctan\dfrac{\sqrt2}{2}$ 时全面积最小
\end{itemize}
